\section{Conclusion}

We present \model{}, a flexible all-atom generator for molecular CSP based on
variational flow matching. \model{} supports multi-component and organometallic
crystals and optional conditioning on molecular templates, stereochemistry, and
space-group information. Controlled studies of architecture, training,
conditioning, inference, and scaling yield a practical recipe based on cacheable
Pairmixer representations, $L^1$ endpoint regression, EDM--Heun sampling, and
balanced single- and pair-track capacity. Inspired by the CCDC CSP blind tests, we introduce a two-track evaluation that separates generator-only coverage from end-to-end performance under a common downstream relaxation and ranking pipeline. Across six benchmarks, \model{} achieves the best matched-budget
coverage, while under the shared downstream pipeline it achieves higher
experimental-form recovery, lower ranks, and faster convergence than the baselines.

\textbf{Limitations and discussion.}
First, our controlled ablations vary design choices one at a time, 
and therefore do not capture interactions among architecture, training,
conditioning, inference, and scaling; the selected recipe is not guaranteed to be
globally optimal. Second, our ablations and model selection primarily use
$k{=}30$ candidates, although the best model may depend on the candidate budget.
Our scaling results show a trade-off between small-budget success and larger-budget
coverage, suggesting that model selection may benefit from larger candidate pools.
Determining an appropriate selection budget therefore remains an important
question for generative CSP.
